\section{引言}

对于诸如大型强子对撞机上希格斯产生这样的高能过程，量子色动力学（QCD）因子化与部分子分布函数（PDF）一直是理论预言的重要工具~\cite{Ellis:1991qj,Lin:2017snn}。
%
但对于涉及观测相对较小横向动量 $Q_\perp$ 的过程，例如 Drell-Yan（DY）产生与半单举深度非弹性散射，则需要一个新的非微扰量——{\it 软函数}（soft function）——来描述在固定 $Q_\perp$ 处不可相消的软胶子辐射物理~\cite{Collins:1981uk,Collins:1984kg,Ji:2004wu,Ji:2004xq}。
%
物理上，DY 中的软函数是这样一个过程的截面：一对高能夸克与反夸克（或胶子）沿相反的光锥方向运动，在湮灭之前辐射出总横向动量为 $Q_\perp$ 的软胶子。尽管在微扰论中计算 $Q_\perp\gg \Lambda_{\rm QCD}$ 区域的软函数已取得很大进展~\cite{Echevarria:2015byo,Li:2016ctv}，但当 $Q_\perp$ 为 ${\cal O}(\Lambda_{\rm QCD})$ 量级时它本质上是非微扰的。从第一性原理计算非微扰的横向动量依赖（TMD）软函数，直到最近才变得可行~\cite{Ji:2019sxk}。


这类计算在格点 QCD 中的主要困难在于：它涉及沿 $n^\pm=\frac{1}{\sqrt 2}(1,\vec 0_\perp,\pm 1)$ 方向（$(t,\perp,z)$ 坐标系）的两条类光 Wilson 线，使得在欧氏空间中直接模拟并不现实。不过，近年来利用大动量有效理论（LaMET）框架计算光锥 PDF 等物理量已取得长足进展~\cite{Ji:2013dva,Ji:2014gla}。
LaMET 的关键观察是：通常由类光场关联子表示的共线夸克与胶子模式~\cite{Collins:2011zzd,Bauer:2000yr,Bauer:2001ct,Bauer:2001yt}，可以在大动量强子态上访问。关于 LaMET 及其在共线 PDF 与其他光锥分布中应用的详细综述可参见文献~\cite{Ji:2020ect,Cichy:2018mum}。最近，本文的部分作者提出，TMD 软函数可以从一个特殊的大动量转移形状因子中提取，该形状因子既可以取自轻介子，也可以取自一对夸克-反夸克色源~\cite{Ji:2019sxk}。
一旦算得该量，DY 及类似过程的 TMD 因子化就可以完全用格点 QCD 可计算的非微扰量来实现~\cite{Ji:2014hxa,Ji:2018hvs,Ebert:2018gzl,Ebert:2019okf,Ji:2019ewn,Vladimirov:2020ofp}。


TMD 软函数通常不是在动量空间、而是在横向坐标空间中借助傅里叶变换变量 $b_\perp $ 来定义和使用的。此外，它还依赖于紫外（UV）重正化标度 $\mu$（通常在维数正规化与最小减除即 $\overline{\rm MS}$ 方案中定义）以及快度调节子 $Y+Y'$
\cite{Collins:2011zzd,Ji:2019sxk}，
\begin{align}
S(b_\perp,\mu,Y+Y')=e^{(Y+Y')K(b_\perp,\mu)}S_I^{-1}(b_\perp, \mu)
\end{align}
其中第一个因子与快度演化有关【由 Collins-Soper（CS）核 $K$ 描述】，第二个因子 $S_I$ 是软贡献中内禀的、与快度无关的部分。快度方案无关的 CS 核 $K$ 已被证明可以通过取两个不同动量下准 TMDPDF 之比来计算~\cite{Ebert:2018gzl,Ebert:2019okf,Ji:2019ewn,Vladimirov:2020ofp,Ebert:2019tvc,Shanahan:2020zxr}。
另一方面，在格点上计算内禀软函数此前从未被尝试过。

本文中，我们在一个 $a=0.098$~fm 的 2+1 味 CLS 系综~\cite{Bruno:2014jqa}上，用若干动量首次给出了内禀软函数 $S_I$ 的格点 QCD 计算，见表~\ref{table:cls}。特别地，我们模拟了大动量轻介子形状因子与准 TMD 波函数（TMDWF），二者之比给出内禀软函数~\cite{Ji:2019sxk}。我们将用 Wilson 圈矩阵元去除准 TMD 波函数中的线性发散。CS 核 $K$ 也可以由准 TMD 波函数对外动量的依赖性计算得到~\cite{Ji:2020ect}，我们将其作为副产品一并算出。我们的结果与 TMDPDF 的 quenched 格点计算结果一致~\cite{Shanahan:2020zxr}。

\begin{figure}[!th]
\begin{center}
\includegraphics[width=0.45\textwidth]{images/01_form_factor.pdf}
\caption{本文计算的赝标介子形状因子 $F$ 示意图。初、末态介子动量大且方向相反。
跃迁``流''由两个局域算符构成，其空间间隔固定为 $b_\perp$。$t_{\rm sep}$
为介子源与汇之间的时间间隔。}\label{fig:formfactor}
\end{center}
\end{figure}
